% Results and discussion: same section order as the methodology chapter.

\chapter{Resultados e discussão}
\label{chap:results}

\placeholder{Apresentar, em um parágrafo, a organização dos resultados, que segue a ordem das seções do \autoref{chap:methodology}.}

\section{Requisitos e arquitetura}
\label{sec:results-architecture}

\placeholder{Apresentar os requisitos definidos e a arquitetura implementada.}

\section{API REST}
\label{sec:results-api}

\placeholder{Apresentar os recursos da API, exemplos de requisição e resposta e a integração com o modelo.}

\section{Interface \textit{web}}
\label{sec:results-web}

\placeholder{Apresentar as telas da interface com figuras.}

\section{Avaliação da plataforma}
\label{sec:results-evaluation}

\placeholder{Apresentar e discutir os resultados da \autoref{tab:evaluation}.}

% Values to be filled in after the tests.
\begin{table}[!htb]
  \centering
  \caption{Resultados da avaliação por atividade}
  \label{tab:evaluation}
  \begin{tabular}{cccc}
    \toprule
    Atividade & Código funcional & Tentativas & Tempo de resposta (s) \\
    \midrule
    A1 & -- & -- & -- \\
    A2 & -- & -- & -- \\
    A3 & -- & -- & -- \\
    \bottomrule
  \end{tabular}
  \fonte{elaborado pelo autor.}
\end{table}
